\inicializarSubseccion{Analisis de la segunda pregunta (Por parte de Viviana)}
\textbf{2. ¿De estos aspectos cuáles considera prioritarios para el diseño del adaptador? selecciona dos opciones.}

\tab De la misma forma que en la pasada en esta pregunta se formaron cuatro conjuntos ya que había cuatro diferentes opciones los cuales los definimos como $E,F,G,H$:

Recordando que $\|U\| = 100$
\begin{itemize}
     \item $E$ = Conservación de temperatura con 69
     \item $F$ = Diseño compacto con 44
     \item $G$ = Fácil de limpiar con 57
     \item $H$ = Que tenga otras cualidades (Add-Ins) con 14
\end{itemize}

\tab Lo cual me implica nuevamente que si o si tengo intersecciones al 182 ser mayor a 100 el cual es mi universo, debido a que de igual forma esta pregunta era opción múltiple, sólo que a diferencia aquí es un poco menor la diferencia.

\tab Con estos datos es fácil visualizar que la moda es $E$, sin embargo, tambien podemos ver que un número importante de personas también escogieron $G$. Esto nos da una idea que el diseño del adapatador tiene que ser fácil de acceder para que este sea fácil de limpiar. 

\tab Esto lo podemos comprobar utilizando la siguiente fórmula de probabilidad de la unión de dos eventos:

\begin{equation*}
    P(E\cup G) = P(E) + P(G) - P(E\cap G)
\end{equation*}

\tab De esta forma podemos ver que alrededor del 82\% de las personas seleccionaron alguna de las dos opciones, lo cual nos da una idea de que ambas son importantes para el diseño del adaptador.